%! TeX program = lualatex
% Two functions for the limit-law examples in lessons/limit-laws.tex: one
% asks for the sum and the difference at 2, the next for the two
% compositions.  Page 1 is the gentle pair, page 2 the one that forces the
% reader to watch which side each limit is approached from, and page 3 the
% exercise, where shifts inside the compositions move the action to x = -1,
% and page 4 the midterm exercise that reads one function at x and at -x.
% Both jump at x = 2, yet both composites still have a limit there:
%
%   lim_{x->2^-} g = 3,   lim_{x->2^+} g = 4    (different, both exist)
%   lim_{x->2^-} f = 0,   lim_{x->2^+} f = 1    (different, both exist)
%
% f is constant (-1) on [2.5, 4.5], an interval holding both of g's
% one-sided limits, so f(g(x)) -> -1 from either side.
% g is constant (3) on [-1, 2), an interval holding both of f's one-sided
% limits, so g(f(x)) -> 3 from either side.
% Neither function is constant, and the sum has no limit at 2: it goes to
% 0 + 3 from the left and 1 + 4 from the right.  That contrast is the point
% of the example -- the sum law fails where the composition survives.
\documentclass[tikz]{standalone}
\input{../typography.tex.preamble}
\input{../tikz.tex.preamble}

\pgfplotsset{composition axis/.style={
    basic curve,
    axis equal image=true,
    % No axis names: reading a composition off these graphs is partly about
    % relabelling the axes -- f's input axis is where g's output lands --
    % so the picture leaves that naming to the reader.  Tick numbers stay;
    % they are what the values are read from.
    xlabel={}, ylabel={},
    width=2.9in, height=2.3in,
    grid=both, minor tick num=1,
    title style={yshift=2pt},
}}

\begin{document}
\begin{tikzpicture}
  \begin{axis}[composition axis, name=fplot,
               xmin=-1.2, xmax=5.8, ymin=-3.2, ymax=2.3,
               title={\(f(x)\)}]
    \addplot[very thick, domain=-1:2]     {x-2};
    \addplot[very thick, domain=2:2.5]    {1-4*(x-2)};
    % the flat stretch that carries both one-sided limits of g
    \addplot[very thick, domain=2.5:4.5]  {-1};
    \addplot[very thick, domain=4.5:5.5]  {-1+2*(x-4.5)};
    \draw[fill=white] (axis cs:2,0) circle (2pt);
    \draw[fill=black] (axis cs:2,1) circle (2pt);
  \end{axis}
  \begin{axis}[composition axis, name=gplot,
               at={(3.4in,0)}, anchor=south west,
               xmin=-2.2, xmax=4.2, ymin=0, ymax=6,
               title={\(g(x)\)}]
    \addplot[very thick, domain=-2:-1] {3+2*(x+1)};
    % the flat stretch that carries both one-sided limits of f
    \addplot[very thick, domain=-1:2]  {3};
    \addplot[very thick, domain=2:4]   {4+(x-2)/2};
    \draw[fill=white] (axis cs:2,3) circle (2pt);
    \draw[fill=black] (axis cs:2,4) circle (2pt);
  \end{axis}
\end{tikzpicture}
% ---------------------------------------------------------------------
% A harder pair, for reading f(g(x)) off the graphs.  Same jumps at 2, but
% now the direction of each approach decides which one-sided value of f to
% read, because f jumps at both of g's one-sided limits:
%
%   x -> 2^-:  g decreases to g(2) = 4, so the argument arrives from ABOVE
%              4, and f(g(x)) -> lim_{t->4^+} f(t) = -1.
%   x -> 2^+:  g rises to 3 from BELOW, so f(g(x)) -> lim_{t->3^-} f = -1.
%
% Both sides give -1, so the limit exists -- but reading f at the dots
% instead, f(4) = f(3) = 0.5, gives the wrong answer.  That is the trap the
% picture is built around.
%
% g is still flat (5) across an interval holding both one-sided limits of f
% at 2, which are 0 and 1, so g(f(x)) -> 5 and that composition is safe.
\begin{tikzpicture}
  \begin{axis}[composition axis, name=fplot,
               xmin=-1.2, xmax=5.8, ymin=-3.2, ymax=1.6,
               title={\(f(x)\)}]
    \addplot[very thick, domain=-1:2]   {x-2};
    \addplot[very thick, domain=2:3]    {1-2*(x-2)};
    % between the two jumps: dips and comes back, so it is constant nowhere
    \addplot[very thick, domain=3:4]    {0.5+2*(x-3)*(x-4)};
    \addplot[very thick, domain=4:5.5]  {-1+(x-4)/2};
    \draw[fill=white] (axis cs:2,0)    circle (2pt);
    \draw[fill=black] (axis cs:2,1)    circle (2pt);
    \draw[fill=white] (axis cs:3,-1)   circle (2pt);
    \draw[fill=black] (axis cs:3,0.5)  circle (2pt);
    \draw[fill=white] (axis cs:4,-1)   circle (2pt);
    \draw[fill=black] (axis cs:4,0.5)  circle (2pt);
  \end{axis}
  \begin{axis}[composition axis, name=gplot,
               at={(3.4in,0)}, anchor=south west,
               xmin=-2.2, xmax=4.2, ymin=1.4, ymax=5.6,
               title={\(g(x)\)}]
    \addplot[very thick, domain=-2:-1]  {5+2*(x+1)};
    % the flat stretch that carries both one-sided limits of f
    \addplot[very thick, domain=-1:1.2] {5};
    % down to g(2) = 4, so x -> 2^- sends the argument to 4 from above
    \addplot[very thick, domain=1.2:2]  {5-(x-1.2)/0.8};
    % up to 3, so x -> 2^+ sends the argument to 3 from below
    \addplot[very thick, domain=2:4]    {3-(x-2)/2};
    \draw[fill=black] (axis cs:2,4) circle (2pt);
    \draw[fill=white] (axis cs:2,3) circle (2pt);
  \end{axis}
\end{tikzpicture}
% ---------------------------------------------------------------------
% Page 3 -- the exercise, where the shifts move the action off the jumps.
% The point of interest is x = -1, and both composites are shifted:
%
%   f(g(x)+1) as x -> -1.  g jumps at -1: it comes DOWN to g(-1) = 2 on the
%     left and UP to 1 on the right, so the argument of f arrives at 3 from
%     above and at 2 from below.  f jumps at both, and the halves that get
%     used are lim_{t->3^+} f = -1 and lim_{t->2^-} f = -1, so the limit is
%     -1.  Reading f at the dots gives f(3) = 0.5 and f(2) = 1 instead.
%
%   g(f(x+1)) as x -> -1.  The shift puts the inner point at 0, where f
%     jumps: f -> -2 from the left and -2.5 from the right.  g is flat (3)
%     across an interval holding both, so the limit is 3.
%
% Shapes are varied from the earlier pages on purpose -- a parabola into
% the jump at 0, lines of different slopes elsewhere -- so the reasoning
% cannot be pattern-matched from the worked example.
\begin{tikzpicture}
  \begin{axis}[composition axis, name=fplot,
               xmin=-1.8, xmax=4.3, ymin=-3.2, ymax=1.6,
               title={\(f(x)\)}]
    \addplot[very thick, domain=-1.5:0] {-2+x^2};
    \addplot[very thick, domain=0:2]    {-2.5+0.75*x};
    \addplot[very thick, domain=2:3]    {1-0.5*(x-2)};
    \addplot[very thick, domain=3:4]    {-1+0.5*(x-3)};
    \draw[fill=white] (axis cs:0,-2)   circle (2pt);
    \draw[fill=black] (axis cs:0,-2.5) circle (2pt);
    \draw[fill=white] (axis cs:2,-1)   circle (2pt);
    \draw[fill=black] (axis cs:2,1)    circle (2pt);
    \draw[fill=black] (axis cs:3,0.5)  circle (2pt);
    \draw[fill=white] (axis cs:3,-1)   circle (2pt);
  \end{axis}
  \begin{axis}[composition axis, name=gplot,
               at={(3.4in,0)}, anchor=south west,
               xmin=-4.3, xmax=3.3, ymin=-1.6, ymax=4.6,
               title={\(g(x)\)}]
    \addplot[very thick, domain=-4:-3]    {3+(-3-x)};
    % the flat stretch that carries both one-sided limits of f at 0
    \addplot[very thick, domain=-3:-1.5]  {3};
    % down to g(-1) = 2, so x -> -1^- sends g(x)+1 to 3 from above
    \addplot[very thick, domain=-1.5:-1]  {3-2*(x+1.5)};
    % up to 1, so x -> -1^+ sends g(x)+1 to 2 from below
    \addplot[very thick, domain=-1:3]     {1-(x+1)/2};
    \draw[fill=black] (axis cs:-1,2) circle (2pt);
    \draw[fill=white] (axis cs:-1,1) circle (2pt);
  \end{axis}
\end{tikzpicture}
% ---------------------------------------------------------------------
% Page 4 -- the Fall 2023 midterm exercise: one function, read at 2 and at
% -2, for lim_{x->2} f(x) + 2 f(-x).  As x -> 2 the reflected argument -x
% approaches -2 from the other side, so the two jumps cancel:
%
%   x -> 2^-:  f(x) -> -4 and f(-x) -> 2, giving -4 + 4 = 0
%   x -> 2^+:  f(x) ->  2 and f(-x) -> -1, giving  2 - 2 = 0
%
% Axis names are kept here, unlike the composition pages: nothing has to be
% relabelled to read this one.
\begin{tikzpicture}
  \begin{axis}[composition axis, xmin=-5, xmax=5, ymin=-5, ymax=5, title={\(f(x)\)}, xtick={-4, -2, 0, 2, 4}, ytick={-4, -2, 0, 2, 4}]
    \addplot[very thick] coordinates {(-5, -4) (-2, -1)};
    \addplot[very thick] coordinates {(-2,  2) ( 2, -4)};
    \addplot[very thick] coordinates {( 2,  2) ( 5,  5)};
    \draw[thick, fill=white] (axis cs:-2,-1) circle[radius=0.15];
    \draw[thick, fill=black] (axis cs:-2, 2) circle[radius=0.15];
    \draw[thick, fill=black] (axis cs: 2,-4) circle[radius=0.15];
    \draw[thick, fill=white] (axis cs: 2, 2) circle[radius=0.15];
  \end{axis}
\end{tikzpicture}
\end{document}
